% ============================================================================
% Fichier cible : ARIEL_LOGMINER_MEMOIRE_FINAL/chapters/chapitre2_etat_art.tex
% Passage à remplacer : de « \section{Corrélation d'événements} » (inclus) à
%   « \section{Mémoire persistante et feedback humain} » (exclu).
% Agent : Agent 3 — État de l'art et positionnement, chantier 2026-11-05.
% Sources : Cuppens et Miège (IEEE S&P 2002) ; Valeur et al. (IEEE TDSC 2004).
% ============================================================================

\section{Corrélation d'alertes}

\subsection{Des alertes élémentaires aux scénarios}

La corrélation ne se réduit pas à placer des alertes proches dans une même fenêtre.
Cuppens et Miège distinguent la gestion de la base d'alertes, le regroupement d'alertes
similaires, leur fusion, puis la corrélation d'étapes compatibles avec un scénario
d'intrusion \cite{cuppens2002alert}. Valeur et al. proposent de même une chaîne de
composants couvrant normalisation, fusion, vérification, corrélation de scénarios,
priorisation et intention ; ils montrent que l'efficacité de chaque composant dépend
fortement du jeu de données \cite{valeur2004correlation}. Ces travaux classiques font de
la corrélation un processus multi-étapes dont les entrées doivent déjà partager une
sémantique minimale.

Trois familles de relations peuvent être distinguées. La similarité rapproche les
doublons ou variantes d'une même alerte. La coïncidence relie des événements proches dans
le temps et le contexte. Enfin, une relation causale peut associer les préconditions et
les conséquences des étapes d'une attaque. Les deux premières réduisent le volume ou
forment des groupes. La troisième prétend reconstruire un scénario et exige un modèle
plus riche des actions, des actifs et de leurs dépendances.

\subsection{Compromis et protocole d'évaluation}

Une règle trop stricte fragmente un incident entre plusieurs groupes ; une règle trop
large fusionne des événements indépendants. La fenêtre temporelle, les clés d'entité et
la stabilité des identifiants déterminent ce compromis. L'évaluation doit donc comparer
le regroupement produit à des incidents de référence, rapporter précision et rappel, et
inclure des baselines plus simples telles que « hôte + fenêtre » ou « fenêtre seule ».

Le corrélateur d'Ariel Logminer appartient à la famille temporelle et contextuelle : il
regroupe les alertes sur une fenêtre fixe et l'égalité de plusieurs clés. Il ne modélise
ni préconditions d'attaque ni causalité. Son corpus synthétique reprend les clés attendues
par la règle, ce qui permet un test fonctionnel mais crée un circuit fermé. Le résultat
doit être interprété au regard de la littérature : précision élevée sur des clés stables,
mais fragmentation des incidents multi-étapes, précisément lorsque des catégories
différentes devraient être reliées. Une validation externe demanderait des incidents
authentifiés, multi-sources et multi-étapes qui n'ont pas servi à concevoir le
corrélateur.

